\documentclass[../../../include/open-logic-section]{subfiles}

\begin{document}

\olfileid[tr]{sth}{ord-arithmetic}{using-addition}
\olsection{Sıra Sayısı Toplamını Kullanmak}

Sıra sayıları üzerindeki toplamayı kullanarak çeşitli kümelerin ranklarını
\olref[spine][rank]{defnsetrank} anlamında açıkça hesaplayabiliriz:

\begin{lem}\ollabel{rankcomputation}
$\setrank{A}=\alpha$ ve $\setrank{B}=\beta$ ise:
\begin{enumerate}
	\item\ollabel{exrankpow} $\setrank{\Pow{A}}=\alpha\ordplus 1$
	\item\ollabel{exrankpair} $\setrank{\{A,B\}}=\max(\alpha,
	\beta)\ordplus 1$
	\item\ollabel{exrankcup} $\setrank{A\cup B}=\max(\alpha,
	\beta)$
	\item\ollabel{exranktuple} $\setrank{\tuple{A,B}}=\max(\alpha,
	\beta)\ordplus 2$
	\item\ollabel{exranktimes} $\setrank{A\times B}\leq\max(\alpha,
	\beta)\ordplus 2$
	\item\ollabel{exrankunion} $\alpha$ boş ya da bir limit sıra sayısı
	olduğunda $\setrank{\bigcup A}=\alpha$; $\alpha=\gamma\ordplus 1$
	olduğunda $\setrank{\bigcup A}=\gamma$.
\end{enumerate}
\end{lem}

\begin{proof}
İspat boyunca \olref[spine][rank]{ranksupstrict} sonucunu tekrar tekrar
kullanıyoruz.

\emph{\olref{exrankpow}.} $x\subseteq A$ ise
$\setrank{x}\leq\setrank{A}$. Dolayısıyla
$\setrank{\Pow{A}}\leq\alpha\ordplus 1$. Özellikle
$A\in\Pow{A}$ olduğundan $\setrank{\Pow{A}}=\alpha\ordplus 1$.

\emph{\olref{exrankpair}.} \olref[spine][rank]{ranksupstrict} gereği.

\emph{\olref{exrankcup}.} \olref[spine][rank]{ranksupstrict} gereği.

\emph{\olref{exranktuple}.} \olref{exrankpair} iki kez uygulanır.

\emph{\olref{exranktimes}.} Her $a\in A$ ve $b\in B$ için
\olref{exranktuple} gereği
$\setrank{\tuple{a,b}}=\max(\setrank{a},\setrank{b})\ordplus 2$ olur.
Şimdi aynı rank-üst-sınırı sonucunu kullanın.

\emph{\olref{exrankunion}.} $\alpha=\gamma\ordplus 1$ ise
$\setrank{c}=\gamma$ olan bir $c\in A$ vardır ve $A$'nın hiçbir
!!{element} daha yüksek ranklı değildir; dolayısıyla
$\setrank{\bigcup A}=\gamma$. $\alpha$ bir limit sıra sayısıysa $A$,
rankları $\alpha$'ya keyfî ölçüde yaklaşan (ama ondan kesin olarak küçük
olan) !!{element}s içerir. Dolayısıyla $\bigcup A$ da rankları $\alpha$'ya
keyfî ölçüde yaklaşan (ama ondan kesin olarak küçük olan) !!{element}s
içerir ve $\setrank{\bigcup A}=\alpha$ olur. Son olarak $\alpha=0$ ise
$A=\emptyset$ ve $\bigcup A=\emptyset$ olduğundan eşitlik yine geçerlidir.
\end{proof}
\noindent
\olref{exranktimes} içinde neden bir eşitlik değil de
\emph{eşitsizlik} bulunduğunu göstermeyi alıştırma olarak bırakıyoruz.

\begin{prob}
$\setrank{A\times B}=\max(\setrank{A},\setrank{B})$ olacak biçimde $A$ ve
$B$ kümeleri verin. Ardından
$\setrank{A\times B}=\max(\setrank{A},\setrank{B})\ordplus 2$ olacak
biçimde $A$ ve $B$ kümeleri verin. Başka ranklar mümkün müdür?
\end{prob}

Artık bir sıra sayısını ``sonlu'' ya da ``sonsuz'' diye betimlemenin ne
anlama gelebileceğine ilişkin birkaç makul kavramın çakıştığını da
gösterebiliriz:

\begin{lem}\ollabel{ordinfinitycharacter}
Her $\alpha$ sıra sayısı için aşağıdakiler denktir:
\begin{enumerate}
	\item\ollabel{ord:notinomega} $\alpha\notin\omega$, yani $\alpha$ bir
	doğal sayı değildir
	\item\ollabel{ord:omegaplus} $\omega\leq\alpha$
	\item\ollabel{ord:oneplus} $1\ordplus\alpha=\alpha$
	%\alpha \approx \alpha \ordplus 1$
	\item\ollabel{ord:plusone} $\alpha\approx\alpha\ordplus 1$, yani
	$\alpha$ ile $\alpha\ordplus 1$ eş güçlüdür
	\item\ollabel{ord:infinite} $\alpha$ Dedekind-sonsuzdur
	\end{enumerate}
\end{lem}
\noindent
Böylece bir sıra sayısının sonlu (ya da sonsuz) olması için gerekenleri
anlamanın ispatlanabilir biçimde denk beş yolunu elde ettik.

\begin{proof}
\emph{\olref{ord:notinomega} $\Rightarrow$ \olref{ord:omegaplus}.}
Üçlü Karşılaştırma gereği.

\emph{\olref{ord:omegaplus} $\Rightarrow$ \olref{ord:oneplus}.}
$\alpha\geq\omega$ sabit olsun. Sonlu Ötesi Tümevarım gereği, bazı limit
sıra sayısı $\beta$ için $\alpha=\beta\ordplus\gamma$ olacak biçimde en
küçük bir $\gamma$ sıra sayısı (muhtemelen $0$) vardır. Şimdi:
\[
	1 \ordplus \alpha =  
	1 \ordplus (\beta \ordplus \gamma) = 
	(1 \ordplus \beta) \ordplus \gamma =  
	\supstrict_{\delta < \beta} (1 \ordplus  \delta) \ordplus  \gamma = 
	\beta \ordplus  \gamma = 
	\alpha.
\]
\emph{\olref{ord:oneplus} $\Rightarrow$ \olref{ord:plusone}.}
Açıkça $f\colon(\alpha\disjointsum 1)\to(1\disjointsum\alpha)$ biçiminde
!!a{bijection} vardır. $1\ordplus\alpha=\alpha$ ise
$g\colon(1\disjointsum\alpha)\to\alpha$ biçiminde bir izomorfizma vardır.
Şimdi $\comp{f}{g}$ bileşkesini ele alın.

\emph{\olref{ord:plusone} $\Rightarrow$ \olref{ord:infinite}.}
$\alpha\approx\alpha\ordplus 1$ ise
$f\colon(\alpha\disjointsum 1)\to\alpha$ biçiminde !!a{bijection} vardır.
Her $\gamma<\alpha$ için $g(\gamma)=f(\gamma,0)$ tanımlayın. Bu
!!{injection}, $f(0,1)\in\alpha\setminus\ran{g}$ olduğundan $\alpha$'nın
Dedekind-sonsuz olduğuna tanıklık eder.

\emph{\olref{ord:infinite} $\Rightarrow$ \olref{ord:notinomega}.}
Bu, \olref[z][infinity-again]{naturalnumbersarentinfinite} sonucudur.
\end{proof}

\end{document}
